% Generated draft: requires mathematical and source review.
\documentclass[11pt]{article}
\usepackage[margin=25mm]{geometry}
\usepackage{fontspec}
\setmainfont{Latin Modern Roman}
\usepackage{amsmath,amssymb,mathtools,unicode-math}
\setmathfont{Latin Modern Math}
\usepackage{enumitem,needspace,xurl,hyperref}
\setlength{\parindent}{0pt}
\setlength{\parskip}{5pt}
\setlength{\emergencystretch}{3em}
\newcommand{\N}{\mathbb{N}}
\newcommand{\Z}{\mathbb{Z}}
\newcommand{\Q}{\mathbb{Q}}
\newcommand{\R}{\mathbb{R}}
\newcommand{\C}{\mathbb{C}}
\newcommand{\F}{\mathbb{F}}
\newcommand{\A}{\mathbb{A}}
\newcommand{\PP}{\mathbb{P}}
\newcommand{\E}{\mathbb{E}}
\newcommand{\eps}{\varepsilon}
\newcommand{\dd}{\,\mathrm{d}}
\begin{document}
\section*{UnsolvedMath ID: 10300023}
\textbf{Source problem code:} catalogue entry
\subsection{Definitions and mathematical statement}
Let \(M\) be a three-manifold, let \(\mathscr F\) be a taut codimension-one foliation, and let \(\Gamma\subset M\) be an embedded graph. In the common closed-manifold convention, tautness means every leaf meets a closed curve transverse to the foliation; manifolds with boundary require an appropriate boundary/transverse-arc convention, which the excerpt does not supply.

First give a useful definition of thin position for \(\Gamma\) relative to \(\mathscr F\). This requires a specified complexity and allowed isotopies, rather than silently importing a global height function from classical thin position: a general foliation is not supplied as a family of levels of such a function. The word ``useful'' has no formal criterion in the excerpt.

Second, when \(\Gamma=\tau^1\) is the one-skeleton of a triangulation \(\tau\), ask whether an isotopy can put the leaves into a decomposition by polyhedral disks of bounded index. A natural ambient-isotopy interpretation moves the triangulation relative to the fixed foliation, or equivalently moves the foliation relative to the fixed triangulation. Components cut out by tetrahedra are a possible meaning of disk pieces, but the source does not define ``polyhedral disk'' or its ``index''.

Thus the intended assertion needs to identify the pieces, the index invariant, and whether the bound is uniform over all leaves and pieces for the given data or universal across triangulations/manifolds. It must also specify the foliation regularity, general position and boundary conventions. No normal-surface, Morse or topological index is selected as a substitute, and boundedness is not assigned a sign convention without the definition.

The supplied summary mentions classical thin position of Gabai and Scharlemann--Thompson but reports no verified foliation-relative theory establishing the requested bounded-index form. It supplies no exact theorem references. The 2002 open question and later assessment remain source attributions pending primary-literature review.
\subsection{Short English statement}
Define thin position of a graph relative to a taut foliation. For a triangulation's one-skeleton, can isotopy make every leaf consist of polyhedral disks with a uniform bounded index? The complexity, disk/index definitions and scope of the bound are not supplied.
\subsection{Sources}
\begin{itemize}
\item Local registry: \url{data/unsolvedmath/problems.json}, permanent UnsolvedMath ID \(10300023\), source code catalogue entry; supplied statement, background, and research summary, updated 2026-08-06.
\item Calegari, \emph{Problems in foliations and laminations of 3-manifolds} (2002), Question 7.9: \url{https://arxiv.org/abs/math/0209081}. Supplied source-TeX extraction accessed 2026-07-29; provenance reports an author-source list of 67 numbered questions.
\item Gabai and Scharlemann--Thompson are named in the summary for classical thin position, without separate full citations or a theorem for this relative question.
\end{itemize}
\end{document}
